\documentclass[10pt]{article}
\usepackage[margin=20mm]{geometry}
\usepackage{amsmath,amssymb,mathtools,fontspec,unicode-math,xurl,hyperref}
\IfFontExistsTF{DejaVu Serif}{\setmainfont{DejaVu Serif}}{\setmainfont{Latin Modern Roman}}
\setmathfont{latinmodern-math.otf}
\setlength{\parindent}{0pt}\setlength{\parskip}{6pt}\sloppy
\title{Initial scientific review of
The global stability of the Kaluza–Klein spacetime}
\author{}
\date{Recovery-0024 M25-0290}
\begin{document}
\maketitle
Cécile Huneau, Annalaura Stingo, and Zoe Wyatt
Journal of the European Mathematical Society; DOI: 10.4171/jems/1663
Reviewed source: arXiv:2307.15267v1, 28 July 2023, 80 pages.\par
\section*{Scope and initial assessment}
This report concerns the exact author manuscript identified above. Its frozen PDF has SHA256\par
`3c7e18884ee2e2c8ddbd4b28049f69e86727af3e0cb3082b6787d9c366799abd`.\par
A publisher version was not available for comparison; equivalence with a version of record is not assumed. The pinned arXiv source was acquired, and the disputed formulas were checked in the TeX and in actual images of the corresponding PDF pages.\par
Initial outcome: Major. Two substantive issues affect the written proof. First, hyperboloidal paths used for sharp estimates leave the regions on which their input bounds have been established. This occurs in the exterior argument and again in the scaling-ray treatment of the nonzero modes. Second, the characteristic estimate in Lemma 4.13 is false as stated; an explicit smooth flat-metric example is given below, together with distinct boundary and mass-hypothesis problems in its application. Six local corrections concern definitions, algebra, time-flux bookkeeping and a source cutoff. A separate question about uniform flux positivity as the mass tends to zero is retained as NeedsEvidence, with its precise limitations.\par
None of the examples below is asserted to disprove the global stability theorem for vacuum Einstein data. The conclusions concern the argument in this particular manuscript. The proposed repairs identify work that would be needed; they do not certify a completed replacement proof. Severity reflects the role of each issue in the core argument. This is an initial scientific review, not independent acceptance. The scientific decision remains pending, and the human-review list is empty.\par
\section*{Coverage, hypotheses and proof structure}
All 80 pages were read as extracted text, including both appendices and the bibliography. Actual source-page images were inspected on pages 1, 3, 7–9, 14–16, 19–20, 22–23, 26–27, 36–40, 46–51, 59–63 and 65–77: 43 pages in total. The TeX was examined at the definitions, characteristic lemma, hyperboloidal operator formulas, oscillator argument and relevant cutoff constructions. The cited Huneau–Stingo manuscript [HS] was read in the relevant part of Section 5, and its pages 28–32 were actually viewed. A complete review of that separate paper is not claimed.\par
The main setting is vacuum Einstein evolution on a manifold with spatial topology \(\mathbb R^3\times\mathbb S^1\), allowing dependence on the compact coordinate. The initial metric and second fundamental form satisfy the constraint equations. The lapse and the other initial spacetime components are chosen to satisfy the wave-coordinate reduction. Theorem 1.2 assumes the two weighted Sobolev tiers in (1.10), including the extra derivative with the additional radial weight, and the separate weighted \(L^2\) bound on \(g_0-g^0\). The tail is prescribed by \[h^0_{\boldsymbol\alpha\boldsymbol\beta}
 =\chi(r/t)\chi(r)\frac{M}{r}\delta_{\boldsymbol\alpha\boldsymbol\beta},
 \qquad h^0_{44}=0.\] The distinction between the positive mass parameter, the size of the remainder, and the smallness needed in the flux argument is relevant to P3.\par
The reduced equation for \(h^1\) has weak-null, classical null, and higher-order terms, together with the differentiated mass-tail source. The exterior argument uses weighted energies for \(Z^{\leq N}h^1\) and for one additional ordinary derivative. Wave gauge improves the \(LT\) components; the improved bounds feed the commutator and weak-null estimates. The exterior energy inequality then supplies the flux through \[\widetilde{\mathcal H}=\{(t-1)^2-r^2=1,\ t\geq2\}.\] The actual exterior slice is \(r\geq\sqrt{t^2-2t}\), while the intended interior hyperboloids are \[\mathcal H_s=\{t^2-r^2=s^2\}\cap\mathcal D^i,
 \qquad
 \mathcal D^i=\{t\geq1+\sqrt{1+r^2}\}.\] Consequently \(t\leq s^2/2\) on \(\mathcal H_s\).\par
The interior proof fixes \(N\geq14\) and \(N_1=N-5\). The full energy allows a half-power loss, the zero-mode energy has a smaller loss, and a restricted number of Lorentz-vector-field derivatives admits the intermediate energy hierarchy. The zero-average component has a circle Poincaré inequality; the zero mode does not. Averaging eliminates pure zero-mode interactions from the zero-average equation, but mixed commutators remain. These distinctions were checked in the energy and derivative allocations, rather than treating all modes as a finite Klein–Gordon system.\par
The argument proceeds through hyperboloidal Sobolev and Hardy bounds, second-derivative estimates for the zero modes, improved \(TU\)-component decay, weak-null estimates and high-order commutators. The ray oscillator estimate yields improved nonzero-mode pointwise bounds. Enhanced zero-mode energies and scalar wave estimates then improve the zero-mode pointwise bound, followed by the mutual refinement in Proposition 4.32 and the energy closure in Section 4.10.\par
The ordering \(\zeta_i\ll\gamma_j\ll\delta_k\) is not, by itself, evidence of circularity. The final closure expressly replaces relevant \(\gamma\)-losses by \(c\sigma\) using Proposition 4.32. For finitely many indices, \(\sigma\) can be selected after the constants in that refinement so that \(c\sigma\) is smaller than the necessary gaps between \(\zeta\)’s. The review does not count the preliminary exponent ordering as a defect. The outstanding problems instead concern the domains and hypotheses needed to obtain those improved estimates.\par
\section*{P1: sharp estimates use paths outside their established domains (Major)}
\subsection*{The exterior hyperboloidal path}
In Proposition 3.3, Step 1 on page 22, the proof selects a time \(\tau_k\in[2^k,2^{k+1})\cap[2,T_0)\), then integrates the hyperboloidal derivative through a point \((t,x)\) until the path meets either \(\tau=\tau_k\) or \(r=2\tau\). The bounds being integrated have been established in \(\mathcal D^e\).\par
Let \(s_*^2=t^2-r^2\) denote the constant along this path. The exterior condition on that same hyperbola is \[r(\tau)^2=\tau^2-s_*^2\geq\tau^2-2\tau
 \quad\Longleftrightarrow\quad s_*^2\leq2\tau.\] Thus a backwards portion of a hyperbola through a point on \(\widetilde{\mathcal H}\) immediately enters \(\mathcal D^i\). For example, take \[t=15,\qquad r^2=195,\qquad s_*^2=30,\qquad \tau_k=9.\] Both times lie in the dyadic interval \([8,16)\). At time \(9\), the path has \(r^2=51\), while the exterior slice requires \(r^2\geq63\). The path has not met \(r=2\tau\), since \(r(\tau)<\tau\). The intermediate interval \(9\leq\tau<15\) is outside the domain of the exterior bounds used in (3.19).\par
This is not a peculiarity of choosing \(9\). Any selected \(\tau_k<16\) precedes some \(t\in[8,16)\); the separator point at that later \(t\) gives the same obstruction. Strictly exterior points sufficiently close to the separator also leave the exterior region before reaching the selected time. Consequently the stated endpoint alternatives do not justify the integration for all points being estimated.\par
A repair must use a path contained in the available region, control an additional boundary contribution, or establish estimates in an enlarged region. Selecting a later averaging window may help in part of a slab, but requires a separate argument near a finite bootstrap endpoint. The current review does not supply such a uniform replacement. The improved \(LT\) decay is used in subsequent commutator and flux estimates, so this is more than a notation correction.\par
\subsection*{The scaling-ray initial term}
For Proposition 4.22, pages 61–62, write \[p_\lambda=\left(\frac{\lambda t}{s},\frac{\lambda x}{s}\right),
 \qquad s^2=t^2-r^2.\] The proof of (4.100) treats \(p_{s_0}\) as lying on the truncated \(\mathcal H_{s_0}\), and uses its boundedness in \(t+r\) and the energy \(E^i(s_0,Z^{\leq2}\mathbf W)\). Direct substitution gives \[(t_\lambda-1)^2-r_\lambda^2-1
 =\lambda\left(\lambda-\frac{2t}{s}\right).\] For positive \(\lambda\), the ray lies in \(\mathcal D^i\) precisely when \(\lambda\geq2t/s\). Its natural entry parameter is therefore \[\lambda_0=\max\{s_0,2t/s\},\] not always \(s_0\).\par
A strictly interior example is \[t=49,\quad s=10,\quad r^2=2301,\quad s_0=\frac{21}{10}.\] Here \(49<s^2/2=50\), but \[t_{s_0}=\frac{1029}{100},
 \qquad
 \max_{\mathcal H_{s_0}}t=\frac{s_0^2}{2}=\frac{441}{200},
 \qquad
 \lambda_0=\frac{49}{5}>s_0.\] The starting point lies well outside \(\mathcal H_{s_0}\). For targets approaching the separator at large \(s\), its time is unbounded; the claim that \(1+t+r=O(1)\) at that starting point is therefore unavailable.\par
There is also a direct functional obstruction to the displayed use of the truncated initial energy. Choose a smooth spacetime bump supported near \(p_{s_0}\), outside a neighbourhood of the compact \(\mathcal H_{s_0}\), and multiply it by \(\cos y\). It has zero circle average. All its derivatives vanish on \(\mathcal H_{s_0}\), so \(E^i(s_0,Z^{\leq2}\mathbf W)=0\), but \(Y_{tx}(s_0)>0\) if the bump is nonzero at the ray point. This disproves a trace estimate of the form (4.100) from that truncated energy alone. It is a test of the claimed functional implication, not an Einstein counterexample.\par
The estimates for \(A_{tx}\) and \(B_{tx}\) likewise require justification on the part of the ray before it enters \(\mathcal D^i\). Merely changing the lower integration limit leaves a boundary term \(Y_{tx}(\lambda_0)\) to estimate with the required directional weight. The directly stated exterior bound for a generic zero-average component is of size \(\epsilon t^{-1+\sigma}\) near the separator. At an entry point with \(k=t/s\), one has \(\lambda_0=2k\), \(t_{\lambda_0}=2k^2\); its contribution to \(\lambda_0^{3/2}\left\|\mathbf W\right\|\) is only \(O(\epsilon k^{-1/2+2\sigma})\). This observation does not exclude a sharper estimate, but it does show why the required \(k^{-3/2}\) factor cannot simply be inserted from that bound.\par
The analogous argument in Huneau–Stingo [Retained latex: \textbackslash{}cite[Section\textasciitilde{}5, pp.\textasciitilde{}28--32]\{HS\}] explicitly extends the relevant estimates to whole hyperboloids before applying the ray argument. In particular, its initial term in (5.19) is used in that enlarged setting. Those estimates concern a different quasilinear system and are not a proof of the missing extension for the present Einstein problem.\par
\paragraph*{Effect and required evidence.}\mbox{}\par
The nonzero-mode pointwise estimate (4.12) and the later refinement in Proposition 4.32 depend on this step. A whole-hyperboloid extension, or a boundary-start ray argument with all required traces and coefficient estimates, is needed. Together with the exterior path problem, this leaves two concrete links in the core bootstrap chain unproved as written. No conclusion that the main theorem itself is false follows.\par
\section*{P2: the characteristic lemma and its application (Major)}
\subsection*{A smooth counterexample to the printed inequality}
Lemma 4.13, equation (4.72), has a left-hand side \(t|(\partial_t-\partial_r)u|\), while its boundary term controls \((\partial_t-\partial_r)(ru)\). These quantities differ by a zeroth-order term: \[r(\partial_t-\partial_r)u
 =(\partial_t-\partial_r)(ru)+u.\] The omitted \(u\) cannot be discarded in the stated general lemma.\par
Take \(g=m\), so \(\pi=0\). Choose a smooth radial function \(\rho\) that vanishes for \(r\leq1/2\) and equals \(1\) for \(r\geq1\), and put \[u(t,x)=\frac{\rho(|x|)}{|x|},\] extended by zero near the origin. This is a globally smooth, time-independent function. Its smooth forcing \(F=\Box u=\Delta u\) is supported in \(r<1\); in particular, no singular source at the origin is being used.\par
Set \(s=5\), \(t=10\), and \(x=11e_1\). This point lies in the printed exterior part of \(\mathcal D_s\), since \[t/2<r<2t-3,\qquad t\leq s^2/2.\] The flat characteristic is \(x(\tau)=(\tau+1)e_1\). For \(2\leq\tau\leq10\), its radius is at least \(3\), so \(F=0\) along the whole integration segment. All angular terms vanish because \(u\) is radial, and all \(\pi\)-terms vanish. On the listed lateral boundary, with its inherited condition \(t\geq2\), one has \(r\geq1\), hence \(ru=1\) and \[(\partial_t-\partial_r)(ru)=0.\] Thus the entire right-hand side of (4.72) is zero. The left-hand side is \[t|(\partial_t-\partial_r)u|(10,11e_1)
 =\frac{10}{121}>0.\] All metric smallness hypotheses are satisfied exactly by \(\pi=0\). This is a counterexample to the literal auxiliary lemma, not to Einstein evolution.\par
\subsection*{Boundary geometry and the mass tail}
There are two further problems independent of that missing term.\par
First, the listed inner boundary retains \(r=t/2\) only for \(t\leq8/3\). With \(\pi=0\), the backwards ray from \((t,r)=(10,6)\), taking \(s=6\) in the printed \(\mathcal D_s\), is \(r(\tau)=\tau-4\). It first reaches \(r=\tau/2\) at \((8,4)\), which is excluded by the stated time restriction. The outer boundary \(r=2\tau-3\) would be reached at \(\tau=-1\), outside the domain. The assertion that every such curve meets the listed boundary is false. Also, the printed condition \(r\leq\sqrt{t^2-s^2}\) selects the region above the terminal hyperboloid, whereas a backward propagation argument using estimates already known up to \(s\) needs a carefully specified region below it.\par
Second, the hypothesis \[|\pi_{LL}|\leq \epsilon\frac{|t-r|}{t+r}\] forces \(\pi_{LL}=0\) on \(r=t\). In the proposed application \(\pi=(H^{\boldsymbol\mu\boldsymbol\nu})^\flat\), the asymptotic tail instead gives \[\pi_{LL}(t,t)=-\frac{2M}{t}
       +O(\epsilon t^{-3/2+2\sigma}).\] For fixed \(M>0\) and sufficiently small \(\sigma<1/4\), the second term cannot cancel the first for all large \(t\). The cited improved \(H^1_{LL}\) estimate does not establish the printed hypothesis for the full \(H\).\par
\paragraph*{Effect and repair status.}\mbox{}\par
The exact transport identity suggests retaining a \(u\)-term and correcting the lateral boundary. A characteristic formulation compatible with the mass tail is also required. Simply adding \(1\) to \(|t-r|\) does not by itself verify all of the new characteristic geometry. Proposition 4.12 and the ensuing good-component weak-null estimates use this lemma. Because a replacement and its application are not proved here, this is classified as a substantive proof defect, despite some locally apparent corrections.\par
\section*{P3: flux positivity and the mass parameter (NeedsEvidence)}
Appendix A uses decay of \(H\) and of \(H_{LL}+2\chi(r/t)\chi(r)M/r\) to compare the curved flux on \(\widetilde{\mathcal H}\) with a positive form involving \[\left(\frac{1}{2(1+r^2)}+\frac{M}{2r}\right)|\partial_tW|^2
       +|\widetilde\nabla W|^2\] where the cutoffs are one. The distinction between smallness relative to a fixed \(M\) and smallness uniform in \(M\) must be made explicit.\par
Here is a test of the auxiliary metric assumptions. Near a large point on the separator, let the inverse metric be diagonal, with \[g^{00}=-1-m+a,\quad g^{ij}=(1-m)\delta^{ij},\quad g^{44}=1,
 \qquad m=M/r,\quad a=\epsilon(1+t+r)^{-3/2}.\] The perturbation \(a\) can be smoothly localized near the chosen point. The mass part is the prescribed \(H^0\), and the only added component is \(H^{00}=a\). Thus \[H_{LL}+2M/r=a,\] which has the required residual decay with \(\delta=1/2\). The general smallness bound for \(H\) is also satisfied with small parameters. The metric remains Lorentzian.\par
Choose the gradient at that point as \[\partial_tW=1,\qquad
 \partial_iW=-\frac{x_i}{t-1},\qquad \partial_yW=0.\] Then \(\widetilde\nabla W=0\), and substituting into (A.6) gives \[e^{\mathrm{curv}}_{\widetilde{\mathcal H}}
 =\frac12\left[
 \frac{1}{1+r^2}
 +\frac{M}{r}\left(1+\frac{r^2}{1+r^2}\right)-a\right].\] The squared norm of the separator normal is the negative of twice this expression. Therefore a negative value also means that the separator is not spacelike for this auxiliary metric.\par
For the concrete values \[\epsilon=10^{-2},\qquad M=10^{-8},\qquad r=10^8,
 \qquad t=1+\sqrt{1+r^2},\] the positive terms inside the brackets sum to less than \(3\times10^{-16}\), while \(a\approx3.5355338529\times10^{-15}\). The flux is approximately \(-1.6177669264\times10^{-15}\). The sign does not rely on floating-point precision: \(1+t+r<3r\) implies \[a>\frac{\epsilon}{6r^{3/2}}
   >\frac{3}{r^2}
   >\frac{1}{1+r^2}
      +\frac{M}{r}\left(1+\frac{r^2}{1+r^2}\right).\]\par
This example shows that the decay assumptions alone do not yield a positivity assertion uniform over independently vanishing \(M\). For \(0<\delta<1\), comparison of a residual \(\epsilon r^{-1-\delta}\) with \(r^{-2}+Mr^{-1}\) involves the size of \[\sup_{r\geq1}\frac{\epsilon r^{1-\delta}}{1+Mr},\] which is of order \(\epsilon M^{\delta-1}\) for small \(M\), up to constants depending on \(\delta\). For each fixed positive \(M\), choosing \(\epsilon\) sufficiently small depending on \(M\) may restore the comparison.\par
The constructed metric is not asserted to satisfy the Einstein equations, wave gauge or the initial constraints. Such structure, or an explicitly allowed dependence of the smallness threshold on \(M\), might resolve the application. The requested evidence is the relevant dependence or an additional argument deriving positivity for the actual data class. This is why the nonlinear consequence is NeedsEvidence, rather than a confirmed counterexample to Theorem 1.2.\par
\section*{Local corrections}
\subsection*{P4: consistent energy domains (Minor)}
On pages 9 and 19 the exterior energy integrates over \(r\geq t-1\), whereas \(\Sigma_t^e\) on page 16 is \(r\geq\sqrt{t^2-2t}\). The assertion (3.5) uses the larger slice. This is not a literal consequence of the narrower energy: at \(t=4\), a smooth spacetime bump supported near \(r=2.9\), inside \[\sqrt{t^2-2t}<r<t-1,\] has zero energy on all the stated narrow slices, but a nonzero gradient norm on \(\Sigma_4^e\). The bump is a functional test, not a nonlinear solution.\par
Define the exterior energy over the actual slices of \(\mathcal D^e\), with its corresponding spacetime region. Likewise, the page 16 truncation \(t\geq r+1\) on \(\mathcal H_s\) gives \(t\leq(s^2+1)/2\), while the separator truncation used on page 9 and in (4.18) gives \(t\leq s^2/2\). These definitions should be made consistent. The correction is local; it does not automatically supply the missing path estimates in P1.\par
\subsection*{P5: the mixed coefficient in (1.18) (Minor)}
With the paper’s null-frame conventions, the correct identity is \[4(t/s)^2\pi^{UV}c^{00}_{UV}
 =\pi_{\underline L\underline L}\frac{s^2}{(t+r)^2}
  +\pi_{LL}\frac{(t+r)^2}{s^2}
  +2\pi_{L\underline L}.\] The last coefficient is \(2\), not \(1\). For \(\pi^{00}=1\) and all other coordinate components zero, each of the three lower null components appearing here is \(1\); the left side minus the printed right side is exactly \(1\). The corrected formula was also checked for arbitrary radial tensor components, not only this test tensor.\par
This correction changes numerical constants in subsequent estimates such as (4.65); it does not force a different decay rate.\par
\subsection*{P6: the wave-operator drift and equation signs (Minor)}
For \(\underline\partial_a=\partial_a+(x_a/t)\partial_t\), the exact identity is \[-\Box_{tx}
 =(s/t)^2\partial_t^2
 +2(x^a/t)\underline\partial_a\partial_t-\underline\partial^a\underline\partial_a
 +\left(\frac3t-\frac{r^2}{t^3}\right)\partial_t.\] The sign in front of \(r^2/t^3\) is positive in (4.62). Applying that printed expression to \(u=t\) gives \(2r^2/t^3\), whereas \(-\Box t=0\). Symbolic differentiation including derivatives of the frame coefficients confirms the corrected identity for a general function.\par
The same sign must be corrected in the definition of \(Q\) in Proposition 4.21 and in the expansion following (4.104). With \[P=\mathcal S^2+2\mathcal S+\frac34,\qquad \mathcal S=t\partial_t+x^a\partial_a,\] the printed \(Q\) differs from \(P-s^2(s/t)^2\partial_t^2\) by \(2s^2r^2t^{-3}\partial_t\).\par
There are associated signs to track when isolating the second time derivative. If \[\Box u+A\,u_{tt}+R[u]=G,\qquad -\Box u=(s/t)^2u_{tt}+L[u],\] then \[((s/t)^2-A)u_{tt}=-L[u]+R[u]-G.\] This is the identity to use in place of the inconsistent signs in (4.63)–(4.64). Also, the mass-tail source has a minus sign in (1.8), but is defined with a plus sign in (3.9) and (4.14).\par
Within \(r<t\), the corrected drift has the same \(O(t^{-1})\) size, and the small-coefficient argument still uses \(|A|\leq(s/t)^2/2\). The estimates mainly take absolute values, so these are local algebraic corrections. They should be carried out before claiming the exact PDE identities; they do not resolve P1.\par
\subsection*{P7: oscillator energy and the mass cutoff (Minor)}
Write the ray equation schematically as \[b\omega''-c\omega_{yy}-d\omega'_y=\mathcal R,
 \qquad c=1+(H^{1,44})^\flat.\] The coefficients are independent of \(y\). For \(E=\int_{\mathbb S^1}(b|\omega'|^2+c|\omega_y|^2)\,\mathrm dy\), integration by parts gives \[E'=b'\left\|\omega'\right\|_2^2+c'\left\|\omega_y\right\|_2^2
       +2\langle\mathcal R,\omega'\rangle.\] The mixed \(d\)-term integrates to zero. The displayed equality on page 61 omits the \(c'\)-term. For the explicit test \(\omega=\lambda\cos y\), \(b=1\), \(c=1+\alpha\lambda\), the omitted quantity is \(\pi\alpha\lambda^2\), which is nonzero for \(\alpha>0\).\par
The definition of \(A_{tx}\) already includes the scaling derivative of \((H^{1,44})^\flat\), so the missing term can be included in the existing Gronwall budget. In the same argument, several occurrences of \(\chi(t/r)\) should be \(\chi(r/t)\), consistently with (1.7) and the definition of \(a(t,x)\) in the proof. These cutoffs are not interchangeable: their supports differ even inside the cone. Restoring the original cutoff and the missing derivative repairs these local identities, subject to the separate domain issue in P1.\par
\subsection*{P8: uniform time-flux bookkeeping (Minor)}
The existence of \(l\in L^1([2,T_0))\) on each finite bootstrap interval does not alone give a bound independent of \(T_0\). To see the distinction, let a cumulative flux be \[F(t)=t^{2\sigma}-2^{2\sigma},\qquad
 l(t)=t^{-2\sigma}F'(t)=\frac{2\sigma}{t}.\] Its total integral is \(2\sigma\log(T_0/2)\). The function is integrable on every finite interval; the issue is uniformity, not finite-interval integrability.\par
What the polynomial cumulative bound does provide is uniform dyadic control. If \(F(t)\leq C\epsilon^2t^{2\sigma}\), with nonnegative flux density \(f=F'\), then for \(l=f/(C\epsilon^2t^{2\sigma})\), \[\int_R^{2R}l(t)\,\mathrm dt\leq2^{2\sigma}.\] Consequently Cauchy–Schwarz yields \[\int_R^{2R}t^{-a}\sqrt{l(t)}\,\mathrm dt
 \leq C_a R^{1/2-a}.\] The dyadic sum converges for \(a>1/2\). This supplies the appropriate mechanism for the weighted integrals in the exterior proof. Parameter margins must be stated accordingly; one may take \(\sigma<1/6\) for a weight \(a=1-3\sigma\), and arrange \(2\sigma<\rho<\kappa/2\) where the \(\rho\)-dependent estimates require it.\par
A final partial dyadic interval can have arbitrarily short length, so its average cannot automatically be bounded by the inverse of the full dyadic scale. A full-window argument or a separate endpoint treatment is needed. These scalar estimates clarify the bookkeeping; they do not certify the path used in Proposition 3.3, whose domain problem was proved in P1.\par
\subsection*{P9: support of the zero-mode source split (Minor)}
Pages 68–70 split the forcing using \(\chi((r+1/2)/t)\), where this cutoff is one on \(B_{1/2}\) and supported in \(B_1\). Its complement is then treated as supported in \(t\leq r+1/2\). This support claim is false. For example, choose an admissible smooth decreasing cutoff with \(\chi(3/4)=1/2\). At \((t,r)=(10,7)\), \[1-\chi((r+1/2)/t)=\frac12,\qquad t-r=3>\frac12.\] Thus the complementary forcing can be nonzero well inside the claimed separating cone.\par
A suitable replacement is a smooth \(\eta(t-r)\) that is zero for \(t-r\leq1/2\) and one for \(t-r\geq1\). The first forcing is then supported where the interior source estimate in Lemma 4.29 applies, and the complementary forcing lies in \(\mathcal D^e\). The transition strip lies in the overlap where the exterior bounds can be used. The Cauchy data should consistently be imposed at \(t=2\), as in the actual proof of Proposition 4.31; the introductory display on page 68 instead uses \(t=0\). This is a local correction of the decomposition.\par
\section*{Independent checks, evidence and limitations}
The calculations were written and executed by the current primary executor. They use Python 3.13.5 and SymPy 1.13.3, with no execution of author code. The script `verification/0290/independent\_checks.py` completed with exit code \(0\); its SHA256 was identical before and after execution:\par
`513e72def5b732ce3fd0819a2a98b84a7b28ea6a01ece05d785dca4f8217dc0f`.\par
The results include exact rational ray geometry, the smooth flat-metric example, the null-frame conversion for a general radial tensor, exact differential-operator residuals, the oscillator energy defect, and an analytic sign bound for the flux witness. Numerical values only illustrate a sign already established by inequalities. No nonlinear PDE simulation or unbounded matrix computation was used.\par
The complete derivations are given in the report; the machine-readable checks and their actual command, exit code, stdout, stderr and script bindings are in `independent\_checks.json` and `independent\_checks\_execution.json`. The additional exterior path calculation is recorded in `exterior\_path\_check.json`. The pinned source archive, extracted TeX and bibliography, source PDF parsing, rendered pages and separate actual-inspection records are included in the support inventory. The source archive was treated as text, not executed.\par
The final report build is recorded separately in `build\_final.json`: it uses the existing approved Tectonic, an isolated dependency cache, untrusted mode and a cached final pass. Compiler, TeX, dependencies and output bindings, real commands and logs are retained there. PDF parsing and actual viewing of every report page have separate records. Rendering alone is not recorded as viewing. Earlier false render flags and the failed direct read of the cited reference’s first text block are preserved; the latter block was subsequently read through the recorded Python process.\par
No defect is inferred from unavailable author executable code or publisher text. The bibliography was read, but every cited result was not independently reconstructed. In particular, the external whole-hyperboloid argument was checked only for its relevance to P1 and is not imported as a proof for the Einstein system. The tentative concern that Lemma 4.28 needs an ordering between \(\mu\) and \(\nu\) is not retained: the null-coordinate majorant uses integrability of the two weights for \(\mu>1,\nu>1\), and no counterexample or missing ordering was established.\par
The elementary witnesses have deliberately limited roles. The bump tests a trace implication; the smooth radial example refutes the printed auxiliary characteristic inequality; and the Lorentzian metric tests uniform positivity under decay assumptions. None is claimed to satisfy the vacuum Einstein constraints and evolution in Theorem 1.2. A repaired exterior/interior sharp-estimate chain, a mass-compatible characteristic lemma, and the precise mass dependence in the flux argument remain necessary before treating the proof as complete.\par
The mechanical inventory and a non-scientific contract-field checklist were the only auxiliary contributions. Controller-supplied UI observations and the runtime self-description are recorded separately in provenance; no independent backend verification is claimed.\par
9 C. Huneau, A. Stingo, and Z. Wyatt, The global stability of the Kaluza–Klein spacetime, arXiv:2307.15267v1, 28 July 2023. Journal identity: DOI 10.4171/jems/1663. The reviewed text is the pinned 80-page author manuscript.\par
C. Huneau and A. Stingo, Global well-posedness for a system of quasilinear wave equations on a product space, arXiv:2110.13982v1, 26 October 2021, Section 5, pp. 28–32. The retrieved 43-page PDF has SHA256 `b540d13095ef0e73a21a0c3448ee6cfba3453da6be47d9d58a39fd3604a25794`.\par
\end{document}
